\chapter[1966 --- Rous et al.]{1966: Peyton Rous, Charles Brenton Huggins}
\label{ch:1966_Rous}

\section*{Main Contribution}

\textit{Discovered tumor-inducing viruses, establishing that viruses could cause cancer and opening new avenues for cancer research.}

\section{Official Citation}

for his discovery of tumour-inducing viruses

\section{Background and Historical Context}

The 1966 Nobel Prize in Physiology or Medicine was awarded to Peyton Rous for his discovery of tumour-inducing viruses and to Charles Brenton Huggins for his discoveries concerning hormonal treatment of prostatic cancer. This prize recognized two related but distinct breakthroughs in cancer research---the discovery that viruses can cause cancer and the demonstration that hormones can be used to treat cancer.

\subsection{The Cancer Problem}

By the mid-20th century, cancer had become one of the leading causes of death in Western countries. However, the causes of cancer were still poorly understood, and treatment options were limited. Surgery and radiation therapy were the primary treatments, but these approaches were often ineffective against advanced cancers. The work of Rous and Huggins would provide new insights into the causes and treatment of cancer.

\subsection{Peyton Rous}

Peyton Rous was born on October 5, 1879, in Baltimore, Maryland. He studied medicine at the Johns Hopkins University, where he earned his medical degree in 1904. After completing his medical training, Rous joined the Rockefeller Institute for Medical Research, where he would spend most of his career.

Rous was a skilled virologist and cancer researcher who was interested in understanding the causes of cancer. His work on tumour-inducing viruses would eventually earn him the Nobel Prize, although the recognition came more than 40 years after his initial discovery.

\subsection{Charles Brenton Huggins}

Charles Brenton Huggins was born on September 22, 1901, in Halifax, Nova Scotia, Canada. He studied medicine at the University of Michigan, where he earned his medical degree in 1924. After completing his medical training, Huggins joined the faculty at the University of Chicago, where he would spend most of his career.

Huggins was a skilled urologist and cancer researcher who was interested in understanding the relationship between hormones and cancer. His work on hormonal treatment of prostatic cancer would earn him the Nobel Prize.

\subsection{The Viral Theory of Cancer}

The idea that viruses could cause cancer had been proposed in the early 20th century, but it was controversial and was not widely accepted. Many researchers believed that cancer was caused by chemical or physical agents rather than by viruses. The work of Rous would provide the first strong evidence for the viral theory of cancer.

\section{The Discovery/Contribution}

The work of Rous and Huggins on cancer involved two related but distinct breakthroughs: the discovery of tumour-inducing viruses and the demonstration that hormones can be used to treat cancer.

\subsection{Rous' Discovery of Tumour-Inducing Viruses}

Peyton Rous' primary contribution was the discovery of the Rous sarcoma virus (RSV), the first tumour-inducing virus to be discovered.

\subsubsection{The Discovery}

In 1911, Rous was studying a type of cancer called sarcoma in chickens. He discovered that the tumour could be transmitted from one chicken to another by injecting a cell-free extract from the tumour into healthy chickens. This result was surprising because it suggested that the tumour was caused by an infectious agent that could pass through filters that would trap bacteria.

Rous demonstrated that the infectious agent was a virus---later named the Rous sarcoma virus (RSV). He showed that the virus could transform normal cells into cancer cells, causing the cells to grow uncontrollably and form tumours.

\subsubsection{Significance of the Discovery}

Rous' discovery was significant because it provided the first strong evidence that viruses could cause cancer. However, the discovery was controversial at the time, and many researchers were skeptical of the idea that viruses could cause cancer. It was not until the 1950s and 1960s, when other tumour-inducing viruses were discovered, that the viral theory of cancer gained widespread acceptance.

\subsubsection{The Oncogene Hypothesis}

Rous' discovery also contributed to the development of the oncogene hypothesis, which proposes that cancer is caused by the activation of specific genes called oncogenes. It was later discovered that the Rous sarcoma virus contains a gene called v-src (viral src) that is responsible for its ability to transform cells. This gene was found to be related to a normal cellular gene called c-src (cellular src) that plays a role in cell growth and division. The discovery of oncogenes was a major breakthrough in cancer research and has led to the development of new treatments for cancer.

\subsection{Huggins' Work on Hormonal Treatment of Prostatic Cancer}

Charles Huggins' primary contribution was the demonstration that hormones can be used to treat cancer, specifically prostatic cancer.

\subsubsection{The Discovery}

Huggins discovered that the growth of prostatic cancer is dependent on the male sex hormone testosterone. He showed that reducing the levels of testosterone in the body---either by surgical removal of the testes (castration) or by the administration of estrogens (female sex hormones)---could cause prostatic cancer to shrink and slow its growth.

\subsubsection{The Experimental Approach}

Huggins' approach was based on his observation that the growth of prostatic cancer cells is stimulated by testosterone. He hypothesized that if testosterone production was reduced, the cancer cells would stop growing and might even die.

Huggins tested this hypothesis in patients with advanced prostatic cancer. He showed that castration or the administration of estrogens could cause significant regression of the tumour, reducing symptoms and improving survival.

\subsubsection{Significance of the Discovery}

Huggins' discovery was significant because it demonstrated that hormones could be used to treat cancer---a concept that was significant at the time. It also provided the first evidence that the growth of some cancers could be controlled by manipulating the hormonal environment of the body.

\subsubsection{Hormonal Therapy for Cancer}

Huggins' work laid the foundation for hormonal therapy, a type of cancer treatment that uses hormones or hormone-blocking agents to control the growth of cancer. Hormonal therapy is now widely used to treat a variety of cancers, including prostatic cancer, breast cancer, and endometrial cancer.

\section{Impact on Modern Medicine}

The work of Rous and Huggins on cancer has had a significant impact on modern medicine. Their discoveries have contributed to the understanding and treatment of cancer and have led to the development of numerous medical technologies and therapies.

\subsection{Viral Oncology}

Rous' discovery of tumour-inducing viruses led to the development of the field of viral oncology---the study of how viruses cause cancer. The discovery of oncogenic viruses has provided important insights into the mechanisms of cancer and has led to the development of vaccines against cancer-causing viruses.

\subsection{HPV Vaccine}

One of a major applications of viral oncology has been the development of vaccines against human papillomavirus (HPV), a virus that causes cervical cancer and other types of cancer. The HPV vaccine, which is based on the principles of viral oncology established by Rous, has been shown to significantly reduce the risk of cervical cancer and other HPV-related cancers.

\subsection{Hormonal Therapy}

Huggins' work on hormonal treatment of prostatic cancer has had a major impact on the treatment of cancer. Hormonal therapy is now widely used to treat a variety of cancers, including prostatic cancer, breast cancer, and endometrial cancer. The development of new hormonal therapies, such as aromatase inhibitors and luteinizing hormone-releasing hormone (LHRH) agonists, has improved the treatment of these cancers.

\subsection{Targeted Therapy}

The understanding of oncogenes that was developed from Rous' work has also contributed to the development of targeted therapy for cancer. Targeted therapy uses drugs that specifically target the molecular abnormalities that cause cancer, such as oncogenes. This approach has been more effective and less toxic than traditional chemotherapy.

\subsection{Cancer Prevention}

The understanding of viral oncology has also contributed to cancer prevention. Vaccines against cancer-causing viruses, such as HPV and hepatitis B virus, can prevent infections that lead to cancer. These vaccines have the potential to prevent millions of cases of cancer worldwide.

\section{Legacy}

The legacy of Peyton Rous and Charles Huggins extends far beyond their Nobel Prize-winning work on cancer. Their discoveries have fundamentally changed our understanding of cancer and have contributed to numerous advances in the diagnosis and treatment of the disease.

\subsection{Foundation for Cancer Research}

The work of Rous and Huggins laid the foundation for modern cancer research. Rous' discovery of tumour-inducing viruses led to the development of viral oncology and the discovery of oncogenes, while Huggins' work on hormonal therapy established the principle that cancer can be treated by manipulating the hormonal environment of the body.

\subsection{Advances in Treatment}

The discoveries of Rous and Huggins have contributed to numerous advances in cancer treatment, including hormonal therapy, targeted therapy, and vaccines against cancer-causing viruses. These advances have improved the lives of millions of cancer patients and continue to drive progress in cancer treatment.

\subsection{The Delayed Recognition of Rous}

One of a notable aspects of Rous' Nobel Prize was the delay in recognition. Rous made his discovery of the Rous sarcoma virus in 1911, but it was not until 1966---55 years later---that he received the Nobel Prize. This delay was due in part to the controversy surrounding the viral theory of cancer, which was not widely accepted until the 1950s and 1960s when other oncogenic viruses were discovered.

Rous' Nobel Prize was also remarkable because he was 78 years old at the time of the award, making him one of the oldest Nobel laureates in Physiology or Medicine. Despite his advanced age, Rous continued to be active in research until his death in 1970.

\subsection{Inspiration for Future Researchers}

The work of Rous and Huggins continues to inspire researchers in cancer biology. Their success in using innovative experiments to answer fundamental questions about cancer demonstrates the power of basic research to improve human health. Their example has motivated generations of scientists to pursue careers in cancer research.

\subsection{Recognition and Honors}

In addition to the Nobel Prize, Rous and Huggins received numerous other honors and awards for their work. Rous received the National Medal of Science in 1966 and was elected to the National Academy of Sciences. Huggins received the National Medal of Science in 1966 and was elected to the National Academy of Sciences.

Their contributions to cancer research and medicine are remembered and celebrated by researchers and clinicians around the world. Their discoveries continue to influence our understanding of cancer and have led to numerous advances in the diagnosis and treatment of the disease, ensuring that their legacy will endure for generations to come.


\section{Modern Developments}

Rous and Huggins' work on tumor viruses has led to modern cancer virology and immunotherapy. Understanding that viruses cause cancer has led to preventive vaccines (HPV vaccine for cervical cancer, Hepatitis B vaccine for liver cancer). Oncolytic virus therapy (T-VEC, approved 2015) uses modified viruses to treat cancer. Understanding of viral oncogenes has revealed cancer signaling pathways, leading to targeted therapies (imatinib for CML).
